\documentclass[a4paper,14pt]{extarticle}

\input{C:/preamble.tex}

\usepackage{tcolorbox}

\begin{document}



	
	Аналогично представим искомые функции через их ограничения на
	подобласти:
	\[
	\widetilde p
	=
	\bigcup_{\alpha=1}^{N_\tau}\widetilde p^{(\alpha)},
	\qquad
	\widetilde{\boldsymbol u}
	=
	\bigcup_{\alpha=1}^{N_\tau}
	\widetilde{\boldsymbol u}^{(\alpha)},
	\]
	где
	\[
	\widetilde p^{(\alpha)}
	=
	\left.\widetilde p\right|_{\tau_\alpha},
	\qquad
	\widetilde{\boldsymbol u}^{(\alpha)}
	=
	\left.\widetilde{\boldsymbol u}\right|_{\tau_\alpha}.
	\]
	
	На данном этапе это лишь запись точного решения по частям; выбор
	конечномерного приближения будет сделан далее.
	
	Первый блок системы можно переписать через ограничения искомой и
	тестовой функций на подобласти:
	\[
	\begin{aligned}
		\nu\sum_{\alpha=1}^{N_\tau}
		\int_{\tau_\alpha}
		\nabla\widetilde v_x^{(\alpha)}
		\cdot\nabla\Phi_{i,\alpha}\,d\tau
		-
		\sum_{\alpha=1}^{N_\tau}
		\int_{\tau_\alpha}
		\widetilde p^{(\alpha)}
		\frac{\partial\Phi_{i,\alpha}}{\partial x}\,d\tau
		&=
		\sum_{\alpha=1}^{N_\tau}
		\int_{\tau_\alpha}
		f_x\Phi_{i,\alpha}\,d\tau,
		\\
		&\hspace{1cm} i=1,\ldots,m_\varphi.
	\end{aligned}
	\]
	
	Обозначим ограничения тестовых функций на подобласти:
	\[
	\Phi_{i,\alpha}
	=
	\left.\varphi_i\right|_{\tau_\alpha},
	\qquad
	\Psi_{i,\alpha}
	=
	\left.\psi_i\right|_{\tau_\alpha},
	\qquad
	\mathcal Q_{i,\alpha}
	=
	\left.q_i\right|_{\tau_\alpha}.
	\]
	
	Тогда система, записанная через ограничения точного решения и
	тестовых функций на подобласти, имеет вид
	\[
	\left\{
	\begin{aligned}
		\nu\sum_{\alpha=1}^{N_\tau}
		\int_{\tau_\alpha}
		\nabla\widetilde v_x^{(\alpha)}
		\cdot\nabla\Phi_{i,\alpha}\,d\tau
		-
		\sum_{\alpha=1}^{N_\tau}
		\int_{\tau_\alpha}
		\widetilde p^{(\alpha)}
		\frac{\partial\Phi_{i,\alpha}}{\partial x}\,d\tau
		&=
		\sum_{\alpha=1}^{N_\tau}
		\int_{\tau_\alpha}
		f_x\Phi_{i,\alpha}\,d\tau,
		&&i=1,\ldots,m_\varphi,
		\\[2mm]
		\nu\sum_{\alpha=1}^{N_\tau}
		\int_{\tau_\alpha}
		\nabla\widetilde v_y^{(\alpha)}
		\cdot\nabla\Psi_{i,\alpha}\,d\tau
		-
		\sum_{\alpha=1}^{N_\tau}
		\int_{\tau_\alpha}
		\widetilde p^{(\alpha)}
		\frac{\partial\Psi_{i,\alpha}}{\partial y}\,d\tau
		&=
		\sum_{\alpha=1}^{N_\tau}
		\int_{\tau_\alpha}
		f_y\Psi_{i,\alpha}\,d\tau,
		&&i=1,\ldots,m_\psi,
		\\[2mm]
		\sum_{\alpha=1}^{N_\tau}
		\int_{\tau_\alpha}
		\mathcal Q_{i,\alpha}\,
		\operatorname{div}
		\widetilde{\boldsymbol u}^{(\alpha)}\,d\tau
		&=0,
		&&i=1,\ldots,m_q.
	\end{aligned}
	\right.
	\]
	
	В каждом слагаемом под интегралом стоят ограничения функций на одну
	подобласть. На данном этапе $\widetilde{\boldsymbol u}$ и
	$\widetilde p$ остаются неизвестными функциями.
	
	\medskip
	
	Для перехода к системе с конечным числом неизвестных будем искать
	приближения решения в виде линейных комбинаций выбранных глобальных
	тестовых функций:
	\[
	\begin{aligned}
		v_{x,h}
		&=
		\sum_{j=1}^{m_\varphi}U_{x,j}\varphi_j,
		\\
		v_{y,h}
		&=
		\sum_{j=1}^{m_\psi}U_{y,j}\psi_j,
		\\
		p_h
		&=
		\sum_{k=1}^{m_q}P_kq_k.
	\end{aligned}
	\]
	
	Здесь $U_{x,j}$, $U_{y,j}$ и $P_k$ --- неизвестные числовые
	коэффициенты. Функции $v_{x,h}$, $v_{y,h}$ и $p_h$ являются
	приближениями точного решения.
	
	На отдельной подобласти эти разложения принимают вид
	\[
	\begin{aligned}
		v_{x,h}^{(\alpha)}
		&=
		\sum_{j=1}^{m_\varphi}
		U_{x,j}\Phi_{j,\alpha},
		\\
		v_{y,h}^{(\alpha)}
		&=
		\sum_{j=1}^{m_\psi}
		U_{y,j}\Psi_{j,\alpha},
		\\
		p_h^{(\alpha)}
		&=
		\sum_{k=1}^{m_q}
		P_k\mathcal Q_{k,\alpha}.
	\end{aligned}
	\]
	
	Коэффициент при каждой глобальной функции один и тот же для всех
	подобластей. Например,
	\[
	U_{x,1}\varphi_1
	=
	U_{x,1}\Phi_{1,1}
	\cup
	U_{x,1}\Phi_{1,2}
	\cup\cdots\cup
	U_{x,1}\Phi_{1,N_\tau}.
	\]
	
	Таким образом, одна глобальная функция вносит вклад в выражение
	решения на каждой подобласти, где её ограничение не равно нулю.
	
	\medskip
	
	Потребуем выполнения слабой формы для приближённого решения и
	выбранных тестовых функций. Первое уравнение для $i$-й тестовой
	функции имеет вид
	\[
	\begin{aligned}
		\nu\sum_{\alpha=1}^{N_\tau}
		\int_{\tau_\alpha}
		\nabla v_{x,h}^{(\alpha)}
		\cdot\nabla\Phi_{i,\alpha}\,d\tau
		-
		\sum_{\alpha=1}^{N_\tau}
		\int_{\tau_\alpha}
		p_h^{(\alpha)}
		\frac{\partial\Phi_{i,\alpha}}{\partial x}\,d\tau
		=
		\sum_{\alpha=1}^{N_\tau}
		\int_{\tau_\alpha}
		f_x\Phi_{i,\alpha}\,d\tau.
	\end{aligned}
	\]
	
	Подставим разложения
	\[
	v_{x,h}^{(\alpha)}
	=
	\sum_{j=1}^{m_\varphi}
	U_{x,j}\Phi_{j,\alpha},
	\qquad
	p_h^{(\alpha)}
	=
	\sum_{k=1}^{m_q}
	P_k\mathcal Q_{k,\alpha}.
	\]
	
	Вынося числовые коэффициенты из интегралов и меняя порядок конечных
	сумм, получаем
	\[
	\begin{aligned}
		&
		\sum_{j=1}^{m_\varphi}
		U_{x,j}
		\left(
		\nu\sum_{\alpha=1}^{N_\tau}
		\int_{\tau_\alpha}
		\nabla\Phi_{j,\alpha}
		\cdot\nabla\Phi_{i,\alpha}\,d\tau
		\right)
		\\
		&\quad-
		\sum_{k=1}^{m_q}
		P_k
		\left(
		\sum_{\alpha=1}^{N_\tau}
		\int_{\tau_\alpha}
		\mathcal Q_{k,\alpha}
		\frac{\partial\Phi_{i,\alpha}}{\partial x}\,d\tau
		\right)
		=
		\sum_{\alpha=1}^{N_\tau}
		\int_{\tau_\alpha}
		f_x\Phi_{i,\alpha}\,d\tau.
	\end{aligned}
	\]
	
	Введём обозначения
	\[
	\begin{aligned}
		(A_x)_{ij}
		&=
		\nu\sum_{\alpha=1}^{N_\tau}
		\int_{\tau_\alpha}
		\nabla\Phi_{j,\alpha}
		\cdot\nabla\Phi_{i,\alpha}\,d\tau,
		\\
		(B_x)_{ki}
		&=
		\sum_{\alpha=1}^{N_\tau}
		\int_{\tau_\alpha}
		\mathcal Q_{k,\alpha}
		\frac{\partial\Phi_{i,\alpha}}{\partial x}\,d\tau,
		\\
		(F_x)_i
		&=
		\sum_{\alpha=1}^{N_\tau}
		\int_{\tau_\alpha}
		f_x\Phi_{i,\alpha}\,d\tau.
	\end{aligned}
	\]
	
	Тогда $i$-е уравнение первого блока записывается как
	\[
	\sum_{j=1}^{m_\varphi}
	(A_x)_{ij}U_{x,j}
	-
	\sum_{k=1}^{m_q}
	(B_x)_{ki}P_k
	=
	(F_x)_i.
	\]
	
	Все уравнения первого блока вместе имеют вид
	\[
	A_x\boldsymbol U_x
	-
	B_x^{\mathsf T}\boldsymbol P
	=
	\boldsymbol F_x.
	\]
	
	Аналогично преобразуя второй блок и уравнение несжимаемости,
	получаем систему
	\[
	\boxed{
		\begin{pmatrix}
			A_x&0&-B_x^{\mathsf T}\\
			0&A_y&-B_y^{\mathsf T}\\
			B_x&B_y&0
		\end{pmatrix}
		\begin{pmatrix}
			\boldsymbol U_x\\
			\boldsymbol U_y\\
			\boldsymbol P
		\end{pmatrix}
		=
		\begin{pmatrix}
			\boldsymbol F_x\\
			\boldsymbol F_y\\
			\boldsymbol 0
		\end{pmatrix}.
	}
	\]
	
	Оставшиеся блоки определяются следующим образом:
	\[
	\begin{aligned}
		(A_y)_{ij}
		&=
		\nu\sum_{\alpha=1}^{N_\tau}
		\int_{\tau_\alpha}
		\nabla\Psi_{j,\alpha}
		\cdot\nabla\Psi_{i,\alpha}\,d\tau,
		\\
		(B_y)_{ki}
		&=
		\sum_{\alpha=1}^{N_\tau}
		\int_{\tau_\alpha}
		\mathcal Q_{k,\alpha}
		\frac{\partial\Psi_{i,\alpha}}{\partial y}\,d\tau,
		\\
		(F_y)_i
		&=
		\sum_{\alpha=1}^{N_\tau}
		\int_{\tau_\alpha}
		f_y\Psi_{i,\alpha}\,d\tau.
	\end{aligned}
	\]
	
	Матрицы $A_x$ и $A_y$ записаны отдельно, поскольку наборы функций
	$\varphi_i$ и $\psi_i$ в общем случае могут различаться.
	
	\medskip
	
	Каждый элемент глобальной матрицы является суммой вкладов отдельных
	подобластей. Например,
	\[
	\begin{aligned}
		(A_x)_{ij}
		={}&
		\underbrace{
			\nu\int_{\tau_1}
			\nabla\Phi_{j,1}\cdot\nabla\Phi_{i,1}\,d\tau
		}_{\text{вклад }\tau_1}
		+\cdots+
		\\
		&+
		\underbrace{
			\nu\int_{\tau_{N_\tau}}
			\nabla\Phi_{j,N_\tau}
			\cdot\nabla\Phi_{i,N_\tau}\,d\tau
		}_{\text{вклад }\tau_{N_\tau}}.
	\end{aligned}
	\]
	
	Таким образом, одна пара глобальных функций $\varphi_i,\varphi_j$
	задаёт один элемент $(A_x)_{ij}$. Для его вычисления интегрирование
	выполняется на каждой подобласти, после чего полученные вклады
	складываются.
	
	Если замыкания соседних подобластей пересекаются по отрезку или
	кривой, это не нарушает разбиения интеграла по площади: общая граница
	имеет двумерную меру нуль. При этом ограничения глобальной функции
	скорости на соседние подобласти должны быть согласованы на общей
	границе; в общем случае это условие понимается в смысле следов.
	
	Для однозначного определения давления дополнительно учитывается
	условие нулевого среднего. Пространства скорости и давления также
	должны быть выбраны так, чтобы обеспечить устойчивость дискретной
	задачи.

\end{document}